% Characteristic-two plane septic case: separate pending independent audit.
\subsection{The exact group for ordinary plane septics in characteristic two}
\label{subsec:septic-exact}

\begin{theorem}\label{thm:septic-stabilizer-monodromy}
For the universal ordinary smooth plane septic in characteristic two,
the geometric and arithmetic monodromy images on
$V=H^1_{\mathrm{et}}(C_{\bar\eta},\mathbf F_2)$ are the full stabilizer
of the distinguished nonzero covector supplied by
$\mathcal O_C(2)$ and the canonical theta characteristic:
\[
 G_{\mathrm{geom}}=G_{\mathrm{arith}}
 =\operatorname{Stab}(\ell)
 \simeq 2^{14}\!:\mathrm{GL}_{14}(\mathbf F_2).
\]
In particular, the rank-fifteen representation is reducible, with its
distinguished invariant hyperplane, and its determinant is trivial.
\end{theorem}

\begin{proof}
We work first over an algebraic closure.  An ordinary genus-fifteen
curve has $2^{15}$ geometric theta characteristics.  Its canonical
theta characteristic $B_C$ is the unique non-effective one.  The
plane theta characteristic $L_C=\mathcal O_C(2)$ has six sections.
The class $v=L_C\otimes B_C^{-1}$ is the distinguished nonzero element
of $V^\vee$, fixed by arithmetic as well as geometric monodromy.
We prove transitivity on the other $2^{15}-2=32766$ elements of
$V^\vee$, then identify the whole stabilizer using a boundary subgroup.

\paragraph{Smooth contact quartics.}
Every nonzero section $s$ of an effective theta characteristic
$\theta_C$ supplies a unique quartic $Q$ with
$Q\cap C=2\operatorname{div}(s)$: the section $s^2$ belongs to
$H^0(C,\omega_C)=H^0(C,\mathcal O_C(4))$, and restriction from plane
quartics is an isomorphism.  If $Q$ is smooth, set
$D=\operatorname{div}(s)$ and
$\theta_Q=\mathcal O_Q(D-3H)$.  The contact reciprocity calculation
\eqref{eq:contact-reciprocity} gives
\[
 \theta_Q^2\simeq\omega_Q,\qquad
 h^0(\theta_C)=1+h^0(\theta_Q).
\]
For every smooth quartic and every theta characteristic on it,
$\theta_Q(3)$ has degree fourteen and exactly twelve sections, by
Riemann--Roch and Serre duality.  Restriction of septic forms to a
quartic is surjective and has kernel
$Q\cdot H^0(\mathbf P^2,\mathcal O(3))$, of dimension ten.

Consider first ordinary smooth quartics with $\theta_Q=B_Q$.
After a purely inseparable Frobenius base change, $B_Q$ is supplied
by the relative Frobenius sequence, so its twelve-dimensional spaces
of sections form a vector bundle.  Squaring their sections and lifting
through restriction gives an irreducible incidence of dimension
\[
 14+(12+10-1)=35,
\]
the dimension of the projective space of septics.  This incidence is
geometrically irreducible: the base is a nonempty open in the
geometrically irreducible quartic parameter space, the Frobenius
base change is a universal homeomorphism, and the section and lifting
spaces have geometrically irreducible affine fibres before
projectivizing.

For all smooth quartics and all theta characteristics, the incidence
dimension is at most thirty-five, since the scheme of roots is
finite over the quartic parameter space and the same twelve-section
calculation applies.  Those with effective $\theta_Q$ have images
of dimension at most thirty-four.  Indeed, reciprocity gives
$h^0(\theta_C)\ge2$, and the open part of
$\mathbf P H^0(C,\theta_C)$ near a smooth contact quartic is a
positive-dimensional fibre of this incidence over $C$.  This covers
the nonordinary smooth quartics as well, since all their geometric
theta characteristics are effective.  Thus a generic septic can
have a smooth contact quartic only through the canonical
non-effective theta characteristic of an ordinary quartic.

We next exclude singular contact quartics for the generic septic,
apart from the squares that give $L_C$.

\paragraph{Nonreduced quartics.}
A nonreduced quartic is either a square of a conic, or has the form
$Q=L^2R$ with $R$ a reduced conic.  A square gives exactly
$\theta_C=\mathcal O_C(2)$.  In the other case, evenness of $Q|C$
implies evenness of $R|C$.
For a fixed smooth conic, the restrictions of septics have binary
degree fourteen.  Those with even zero divisor are squares and form
an eight-dimensional space; the restriction kernel has dimension
$h^0(\mathbf P^2,\mathcal O(5))=21$.
Allowing the conic to vary gives dimension
$5+8+21-1=33$ for the image in the septic parameter space.
For a reduced pair of lines, each restricted binary form has odd
degree seven.  Its sole possible odd-order zero must be the node,
so each restriction is the fixed linear node factor times the
square of a cubic.  This leaves four parameters on each line;
both node values vanish, so their combined dimension is eight.
The conic parameter dimension is four, giving the smaller bound
$4+8+21-1=32$.  Consequently all non-square nonreduced contact
quartics occur over a proper subset of the septic parameter space.

\paragraph{Reduced reducible quartics.}
We give upper bounds that allow all coincidences and higher contacts
among components.  For a fixed reduced quartic, normalize each
component.  Away from the preimages of singular points, the divisor
of the restricted septic is even.  For each choice of the finite
odd-parity divisor $\epsilon$ on these preimages, its restriction
is a fixed section with divisor $\epsilon$ times a square.
The square-root line bundles form a finite set.  On a normalization
component of degree $a$ and genus $h$, the root has degree
$(7a-|\epsilon_i|)/2$ and, in the cases below, vanishing $H^1$;
hence its section dimension is
$(7a-|\epsilon_i|)/2+1-h$.
These degrees exceed $2h-2$: the only positive-genus component is
a smooth cubic, where the root degree is at least nine.
On a line the parity divisor has at most three points, and on a
conic at most four, so the genus-zero assertion is immediate.

For a smooth cubic and a line, the base dimension is eleven.
Both components have odd total intersection degree, so
$|\epsilon|\ge2$.  The sum of root dimensions is at most
$14+2-1-1=14$.  Adding the ten-dimensional restriction kernel and
projectivizing gives at most $11+14+10-1=34$.
If the cubic is irreducibly singular, its base dimension decreases
by at least one while the normalization genus decreases by one;
the same upper bound thirty-four follows.  Cases where the cubic
itself is reducible are included below.

For two smooth conics the base dimension is ten and the unrestricted
sum of root dimensions is sixteen.  If $\epsilon\ne0$, its total
degree is at least two, giving at most fifteen.  If
$\epsilon=0$, choose any intersection point of the two conics.
Descent of the septic restriction imposes equality of its two values,
and, after taking square roots over the algebraically closed field,
this is one independent linear condition on the root sections.
Independence follows from global generation on both normalization
components.  Thus the root dimension is again at most fifteen,
and the total bound is $10+15+10-1=34$.

For a smooth conic and two distinct lines the base dimension is
nine.  The two odd-degree line restrictions force
$|\epsilon|\ge2$, so the sum of root dimensions is at most
$14+3-1=16$.  The total is at most $9+16+10-1=34$.
For four distinct lines the base dimension is eight.  Each line
forces an odd parity divisor, so $|\epsilon|\ge4$ and the sum of
root dimensions is at most $14+4-2=16$; the total is at most
$8+16+10-1=33$.  No transversality assumption was used in these
upper bounds.

\paragraph{Integral singular quartics: normalization parameter spaces.}
An integral singular quartic has normalization of genus
$h\in\{0,1,2\}$.  We need two parameter bounds, and establish them
directly in characteristic two instead of using a Severi-variety
dimension assertion:
\begin{align*}
 \dim\{Q:\text{normalization genus }h\}&\le11+h,\\
 \dim\{Q:\text{normalization genus }h,
          \text{all singularities unibranch}\}&\le10+h.
\end{align*}
For $h=0$, a normalization map is given by a base-point-free
three-dimensional subspace of $H^0(\mathbf P^1,\mathcal O(4))$,
together with a projective frame.  The parameter dimension is
$\dim\operatorname{Gr}(3,5)+\dim\operatorname{PGL}_3
 -\dim\operatorname{PGL}_2=6+8-3=11$.
For $h=1$, temporarily mark an origin on the normalization.  Its
moduli have dimension one, its degree-four Picard variety dimension
one, the three-dimensional subspace of its four-dimensional section
space varies in dimension three, and a projective frame contributes
eight.  Forgetting the arbitrary origin subtracts one, giving
$1+1+3+8-1=12$.
For $h=2$, smooth-curve moduli have dimension three, the degree-four
Picard variety dimension two, and every degree-four line bundle has
three sections, giving $3+2+8=13$.
The dimension of genus-two moduli follows directly from curve
deformations: $H^0(T)=0$, $h^1(T)=3$ by Riemann--Roch, and $H^2(T)=0$.
The images under consideration are birational normalization maps,
so quotienting by reparametrization introduces no positive-dimensional
stabilizer.  These constructions cover all integral quartics of the
indicated normalization genus.

If all singularities are unibranch, the normalization map is
nonimmersive somewhere.  Otherwise at the preimage of each singular
point, one plane coordinate would be a uniformizer, making that
branch smooth.  For $h=0$ or $1$, and a degree-four line bundle $A$,
the map
$H^0(A)\to H^0(A|_{2P})$ is onto for every $P$, by Serre duality.
Among three-dimensional subspaces of $H^0(A)$, failure to have rank
two under this map has codimension two: for $h=1$ the rank-one locus
has dimension one in $\operatorname{Gr}(3,4)$ of dimension three;
for $h=0$ it has dimension four in
$\operatorname{Gr}(3,5)$ of dimension six.
Allowing $P$ to vary restores at most one dimension, so the
nonimmersive locus has codimension at least one in the parameter
space above.  This argument is invariant under reparametrization,
and the arbitrary elliptic origin remains a one-dimensional fibre.
For $h=2$, the system is the complete system $|A|$.
Failure to separate $2P$ is possible exactly when
$H^1(A(-2P))\ne0$, or $A\simeq\omega(2P)$.
As $P$ varies this is a subset of dimension at most one in the
two-dimensional degree-four Picard variety.  This remains a
dimension bound in characteristic two even if the map
$P\mapsto\omega(2P)$ is inseparable.  Thus the second bound follows
for all three genera.

Return to the restricted septic on such a quartic.  Its pullback
to the normalization has degree twenty-eight.  The odd-parity
divisor $\epsilon$ is supported above the finitely many singular
points.  There are at most six such branches, since the arithmetic
genus defect is at most three and the number of branches is at
most twice that defect.  A root after removing $\epsilon$ therefore
has degree at least eleven and vanishing $H^1$.  Its section
dimension is exactly
\[
 15-h-\frac{|\epsilon|}{2}.
\]
If $\epsilon\ne0$, then $|\epsilon|\ge2$; together with the
parameter bound $11+h$ this gives incidence dimension at most
$(11+h)+(14-h)+10-1=34$.
If $\epsilon=0$ and a singular point has two different branches,
descent imposes equality of the values at two distinct preimages.
The root line bundle separates these points, again by its degree,
so this is one independent linear condition after taking square
roots; the same bound thirty-four follows.
If all singularities are unibranch, the improved parameter bound
$10+h$ instead gives
$(10+h)+(15-h)+10-1=34$ without using any descent condition.
This excludes every integral singular contact quartic over a
generic septic.

\paragraph{Transitivity.}
Outside the proper closed subsets just bounded, every effective
theta characteristic other than $L_C$ has a smooth contact quartic,
and that contact quartic arises from its canonical non-effective
theta characteristic.  All $32766$ such theta characteristics
therefore occur in the geometrically irreducible canonical incidence
of dimension thirty-five.  This proves dominance.  Its generic
fibre is finite; in particular each of these theta characteristics
has exactly one section up to scalar, since otherwise its projective
section space would give a positive-dimensional fibre.  Consequently
the generic geometric fibre has precisely $32766$ points.
Irreducibility gives geometric Galois transitivity on these points;
a purely inseparable part of the incidence does not affect this
conclusion.  Hence $G_{\mathrm{geom}}$ is transitive on
$V^\vee\setminus\{0,v\}$.

\paragraph{The full quotient on the fourteen-dimensional space.}
Let $W=V^\vee/\langle v\rangle$, and let $H$ be the image on $W$.
Transitivity just proved implies transitivity of $H$ on
$W\setminus\{0\}$, of size $2^{14}-1$.
Degenerate to an ordinary sextic and a transverse line meeting it
in six points.  The rank-fifteen local system is still defined on
this nodal septic because Frobenius is bijective; its graph part
has rank five and its curve quotient has rank ten.  The projection
from this boundary family to the ordinary sextic parameter space
has geometrically connected fibres, so the quotient monodromy is
the full $\mathrm{GL}_{10}(2)$ established above.
In particular $11\mid|G_{\mathrm{geom}}|$.
The kernel of $\operatorname{Stab}(v)\to\mathrm{GL}(W)$ is a
two-group; therefore $11\mid|H|$.

Choose in the sextic quotient an order-thirty-one element acting
as a primitive five-dimensional block plus an identity of rank five.
An order-thirty-one lift exists by the Sylow argument used for
sextics: $\mathrm{GL}_{15}(2)$ has no element of order $31^2$, since
$\operatorname{ord}_{31^2}(2)=155$.
Its action on the graph part is trivial, since the node action
factors through $S_6$.  On $V^\vee$ it therefore has characteristic
polynomial $(T-1)^{10}f(T)$ for a primitive irreducible
degree-five factor $f$.  Quotienting by the fixed $v$ gives on $W$
\[
 (T-1)^9 f(T).
\]
This polynomial is not reciprocal, and the fixed-space dimension
is nine.

Hering's classification, in the opened form of
Churikov--Vasil'ev--Zvezdina (2022), Lemma~3.3, gives, in dimension
fourteen over $\mathbf F_2$, the alternatives
$\mathrm{GL}_{14}(2)$ or containment in
\[
 N(\mathrm{Sp}_{14}(2)),\quad
 \Gamma\mathrm L_7(4),\quad
 \Gamma\mathrm L_2(128),\quad
 \Gamma\mathrm L_1(16384).
\]
The symplectic alternative forces a reciprocal polynomial and is
excluded.  In $\Gamma\mathrm L_7(4)$ an order-thirty-one element
is $\mathbf F_4$-linear, so its fixed dimension over $\mathbf F_2$
is even, excluding nine.  The orders of the last two groups are
$1864254336$ and $229362$, respectively, neither divisible by eleven.
Thus $H=\mathrm{GL}_{14}(2)$.

\paragraph{The unipotent kernel and arithmetic equality.}
The intersection of $G_{\mathrm{geom}}$ with the kernel
$2^{14}$ of $\operatorname{Stab}(v)\to\mathrm{GL}(W)$ is an invariant
subspace for the full natural dual $\mathrm{GL}_{14}(2)$ module.
It is therefore zero or the whole kernel.  In the former case,
$G_{\mathrm{geom}}$ is a complement.  The elementary transvection
calculation $H^1(\mathrm{GL}_m(2),\mathbf F_2^m)=0$ for $m\ge4$,
proved in the quintic section (and equally valid for the dual),
shows that every such complement is conjugate to the standard one.
It has two orbits of size $2^{14}-1$ on
$V^\vee\setminus\{0,v\}$, contradicting transitivity.
Hence the whole kernel occurs, and the geometric group is the
full stabilizer.  Arithmetic monodromy also fixes $v$ and contains
the geometric group, so the two groups are equal.
\end{proof}

\paragraph{Finite certificate and scope.}
The retained program \texttt{sextic\_septic\_group\_checks.py}
checks the order of $2$ modulo $31^2$, the two distinct inverse
Frobenius orbits of primitive thirty-first roots, the semilinear
group orders and their missing prime divisors, and the two orbit
sizes used above.  It does not recompute Hering's classification
or replace the incidence argument by experimental Frobenius data.
This theorem settles degree seven in characteristic two only;
it does not assert the same stabilizer in every odd degree.
